\documentclass{article}
\usepackage{graphicx} % Required for inserting images
\usepackage{amsmath,amssymb,amsthm}
\usepackage[a4paper,margin=1in]{geometry}
\title{PDerivation Of Q(x) and Test}
\author{Bishnu Prasad Swain}
\date{June 2026}

\begin{document}
\title{Derivation Of Q(x) and Test}
\author{Bishnu Prasad Swain}
\date{June 2026}
\maketitle

\bigskip

\section*{Part 1 -- Explicit factorization of \(P(x)\)}

\subsection*{Section 1: Define the variables}

Let
\begin{equation}
A=\cos\alpha,\qquad B=\cos\beta,\qquad C=\cos(\alpha-\beta),\qquad Y=A+B=\cos\alpha+\cos\beta.
\tag{1}
\end{equation}

From the paper,
\begin{equation}
c=\cos\frac{\alpha-\beta}{2}
\quad\Longrightarrow\quad
c^{2}=\frac{1+C}{2}.
\tag{2}
\end{equation}

\subsection*{Section 2: Expand \(P(x)\)}

First, expand the product:
\begin{equation}
(x^{2}+1-2Ax)(x^{2}+1-2Bx)
= x^{4}-2Yx^{3}+(2+4AB)x^{2}-2Yx+1.
\tag{3}
\end{equation}

Multiply by \(c^{2}\):
\begin{equation}
c^{2}x^{4}-2c^{2}Yx^{3}+c^{2}(2+4AB)x^{2}-2c^{2}Yx+c^{2}.
\tag{4}
\end{equation}

Now expand the squared term:
\begin{equation}
\bigl(x^{2}-Yx+C\bigr)^{2}
= x^{4}-2Yx^{3}+(Y^{2}+2C)x^{2}-2YCx+C^{2}.
\tag{5}
\end{equation}

Subtract (5) from (4) to obtain \(P(x)\):
\begin{equation}
\begin{aligned}
P(x)={}&(c^{2}-1)x^{4}
+2Y(1-c^{2})x^{3} \\
&+\bigl[c^{2}(2+4AB)-Y^{2}-2C\bigr]x^{2} \\
&+2Y(C-c^{2})x
+(c^{2}-C^{2}).
\end{aligned}
\tag{6}
\end{equation}

Using (2), simplify the coefficients:

\begin{itemize}
\item \(x^{4}\): \(\displaystyle c^{2}-1=\frac{C-1}{2}\).
\item \(x^{3}\): \(\displaystyle 2Y(1-c^{2})=Y(1-C)\).
\item \(x\): \(\displaystyle 2Y(C-c^{2})=Y(C-1)=-Y(1-C)\).
\item Constant:
\[
c^{2}-C^{2}
=\frac{1+C}{2}-C^{2}
=\frac{2C^{2}-C-1}{2}
=\frac{(2C+1)(C-1)}{2}.
\]
\end{itemize}

So
\begin{equation}
\begin{aligned}
P(x)={}&\frac{C-1}{2}x^{4}
+Y(1-C)x^{3} \\
&+\bigl[c^{2}(2+4AB)-Y^{2}-2C\bigr]x^{2} \\
&-Y(1-C)x+\frac{(2C+1)(C-1)}{2}.
\end{aligned}
\tag{7}
\end{equation}

\subsection*{Section 3: Divide by \((x^{2}-1)\)}

We seek a quadratic
\[
Q(x)=q_{2}x^{2}+q_{1}x+q_{0}
\]
such that
\[
P(x)=(x^{2}-1)Q(x).
\]

Expanding the right-hand side:
\begin{equation}
(x^{2}-1)(q_{2}x^{2}+q_{1}x+q_{0})
= q_{2}x^{4}+q_{1}x^{3}+(q_{0}-q_{2})x^{2}-q_{1}x-q_{0}.
\tag{8}
\end{equation}

Equate the coefficients of (7) and (8):

\begin{itemize}
\item \textbf{\(x^{4}\)}:
\[
q_{2}=\frac{C-1}{2}.
\]

\item \textbf{\(x^{3}\)}:
\[
q_{1}=Y(1-C).
\]

\item \textbf{\(x\)}:
\[
-q_{1}=-Y(1-C),
\]
which is consistent.

\item \textbf{Constant}:
\[
-q_{0}=\frac{(2C+1)(C-1)}{2}
\]
\[
\Longrightarrow
q_{0}
=
-\frac{(2C+1)(C-1)}{2}
=
\frac{(2C+1)(1-C)}{2}.
\]

\item \textbf{\(x^{2}\)}: The identity
\[
q_{0}-q_{2}
=
c^{2}(2+4AB)-Y^{2}-2C
\]
holds by standard trigonometric relations.
\end{itemize}

Thus the quotient is completely determined.

\subsection*{Section 4: Write \(Q(x)\) in compact trigonometric form}

Using the values found:
\[
q_{2}=\frac{C-1}{2},
\qquad
q_{1}=Y(1-C),
\qquad
q_{0}=\frac{(2C+1)(1-C)}{2}.
\]

Therefore
\[
Q(x)
=
\frac{C-1}{2}x^{2}
+
Y(1-C)x
+
\frac{(2C+1)(1-C)}{2}.
\]

Factor out the common factor \(\displaystyle \frac{C-1}{2}\) (note that \(1-C=-(C-1)\)):
\begin{equation}
Q(x)
=
\frac{C-1}{2}x^{2}
-Y(C-1)x
+\frac{(2C+1)(C-1)}{2}
=
\frac{C-1}{2}
\Bigl[x^{2}-2Yx+(2C+1)\Bigr].
\tag{9}
\end{equation}

\subsection*{Section 5: Final factorized form of \(P(x)\)}

Substitute back \(C=\cos(\alpha-\beta)\) and \(Y=\cos\alpha+\cos\beta\):

\begin{equation}
\boxed{
P(x)=(x^{2}-1)\,
\frac{\cos(\alpha-\beta)-1}{2}
\left[
x^{2}-2(\cos\alpha+\cos\beta)x
+2\cos(\alpha-\beta)+1
\right].
}
\tag{10}
\end{equation}

Equivalently, expanding the quotient explicitly:
\begin{equation}
\boxed{
Q(x)=
\frac{\cos(\alpha-\beta)-1}{2}x^{2}
+(\cos\alpha+\cos\beta)\bigl(1-\cos(\alpha-\beta)\bigr)x
+\cos^{2}(\alpha-\beta)-\frac{1+\cos(\alpha-\beta)}{2}.
}
\tag{11}
\end{equation}

\subsection*{Section 6: Where is \(c\) in this \(Q(x)\)?}

The constant \(c\) does not appear as an external multiplier; instead it is embedded in the coefficients because

\[
c=\cos\frac{\alpha-\beta}{2}
\quad\Longrightarrow\quad
c^{2}=\frac{1+\cos(\alpha-\beta)}{2}.
\]

Thus the constant term in (11) becomes

\[
\cos^{2}(\alpha-\beta)-c^{2}.
\]

So \(c\) is inside \(Q(x)\), not outside. This completes the derivation.

\subsection*{Section 7: A simple sine-only version of \(Q(x)\)}

Using the standard identities
\[
\cos\theta = 1 - 2\sin^2\frac{\theta}{2},
\]
the quadratic above simplifies to the beautiful and compact form stated in Lemma 2.3 of the paper:

\begin{equation}
\boxed{
Q(x)=
4\sin^{2}\frac{\alpha}{2}\sin^{2}\frac{\beta}{2}\,(x-1)^2
+
4\sin^{2}\frac{\alpha-\beta}{2}\,x.
}
\tag{12}
\end{equation}

This expression is entirely in terms of sines, has a perfect square as the first term, and a positive coefficient for \(x\) in the second term (since \(\alpha \neq \beta\)). It also makes the presence of \(c\) even clearer, since
\[
\sin^{2}\frac{\alpha-\beta}{2}=1-c^{2}.
\]

\bigskip

\section*{Part 2 -- Positivity of \(Q(x)\) for \(x>0\)}

From the sine-only expression (12), expand the square to rewrite \(Q(x)\) in standard quadratic form:

\begin{equation}
Q(x) = A x^{2} + (B-2A)x + A,
\tag{13}
\end{equation}

where

\begin{equation}
A = 4\sin^{2}\frac{\alpha}{2}\sin^{2}\frac{\beta}{2} > 0,
\tag{14}
\end{equation}

\begin{equation}
B = 4\sin^{2}\frac{\alpha-\beta}{2} > 0
\qquad (\alpha \neq \beta).
\tag{15}
\end{equation}

The discriminant of \(Q(x)\) is

\begin{equation}
\Delta = (B-2A)^2 - 4A^2 = B^2 - 4AB = B(B-4A).
\tag{16}
\end{equation}

We must show \(Q(x) > 0\) for every \(x > 0\).

\subsection*{Case 1: \(\Delta < 0\) (i.e. \(B-4A < 0\))}

When \(\Delta < 0\), the quadratic has no real roots. Because the leading coefficient \(A > 0\), the parabola opens upward and lies entirely above the \(x\)-axis. Hence

\begin{equation}
Q(x) > 0
\quad \text{for all real } x.
\tag{17}
\end{equation}

In particular, \(Q(x) > 0\) for all \(x > 0\).

\subsection*{Case 2: \(\Delta \ge 0\) (i.e. \(B-4A \ge 0\))}

Here \(Q(x)\) has two real roots (possibly equal). Denote them by \(r_1\) and \(r_2\).

\paragraph{Step 1: Vieta's formulas}

\begin{equation}
r_1 + r_2
=
-\frac{B-2A}{A}
=
2-\frac{B}{A}.
\tag{18}
\end{equation}

\begin{equation}
r_1r_2=\frac{A}{A}=1.
\tag{19}
\end{equation}

\paragraph{Step 2: Sign of the sum}

From \(B-4A \ge 0\) we obtain
\[
\frac{B}{A}\ge4.
\]

Therefore

\begin{equation}
r_1+r_2
=
2-\frac{B}{A}
\le
2-4
=
-2<0.
\tag{20}
\end{equation}

So the sum is strictly negative.

\paragraph{Step 3: Sign of the product}

\begin{equation}
r_1r_2=1>0.
\tag{21}
\end{equation}

Hence the two roots have the same sign (both positive or both negative).

\paragraph{Step 4: Determine the common sign}

If both roots were positive, their sum would be positive. But we have shown the sum is negative. Therefore both roots cannot be positive. Consequently,

\begin{equation}
r_1<0,
\qquad
r_2<0.
\tag{22}
\end{equation}

(If \(\Delta=0\), the double root is also negative because the sum is negative and the product positive.)

\subsection*{Conclusion of the two cases}

\begin{itemize}
\item \textbf{Case 1} (\(\Delta<0\)): \(Q(x)>0\) for all real \(x\), hence for all \(x>0\).

\item \textbf{Case 2} (\(\Delta\ge0\)): The two real roots are both negative. Since the quadratic opens upward (\(A>0\)), it is positive for every \(x\) greater than the larger root. The larger root is negative, so for any \(x>0\) we have \(Q(x)>0\).
\end{itemize}

Thus in every possible situation,

\begin{equation}
Q(x) > 0
\quad \text{for all } x > 0.
\tag{23}
\end{equation}

\subsection*{Final deduction for \(P(x)\) and the original equation}

Recall that \(P(x)=(x^2-1)Q(x)\). For any \(x>0\),

\begin{equation}
P(x)=0
\;\Longleftrightarrow\;
(x^2-1)Q(x)=0.
\tag{24}
\end{equation}

Because of (23), \(Q(x)\neq0\) for all \(x>0\). Hence the factor \(Q(x)\) never vanishes on the positive real line. Therefore

\begin{equation}
P(x)=0
\;\Longleftrightarrow\;
x^2-1=0.
\tag{25}
\end{equation}

The equation \(x^2-1=0\) has exactly two real solutions:
\[
x=1
\quad\text{and}\quad
x=-1.
\]

Only \(x=1\) lies in the domain \(x>0\).

Consequently, \textbf{the only positive root of \(P(x)=0\) is \(x=1\)}.

The original unsquared equation is equivalent to \(P(x)=0\) together with the sign condition that both sides are positive. At \(x=1\) the left-hand side equals \(c>0\), so the sign condition is satisfied. If there were any other positive \(x_0\neq1\) satisfying the unsquared equation, it would also satisfy \(P(x_0)=0\), contradicting the uniqueness just proved. Hence \(x=1\) is the \textbf{unique positive solution} of the original geometric condition.

This completes the rigorous proof.

\end{document}

